% ============================================================
%  CHAPTER 19: REFLECTION OF LIGHT
%  The Physics Codex: A Complete University Foundation
%  Korede Samuel Omotosho (Falcon)
%  Aurikrex Learning Systems, 2026
%
%  File: chapters/part3/ch19_reflection.tex
% ============================================================

\chapter{Reflection of Light}
\label{ch:reflection}
\minitoc
\newpage

\chapterquote{We are all mirrors and lenses for one
another --- every person you meet reflects something
back to you, and focuses your attention on something
new.}{Anonymous}

\begin{objectivebox}
By the end of this chapter, you should be able to:
\begin{enumerate}[noitemsep]
  \item State and apply the laws of reflection
  \item Describe the properties of images formed
  in plane mirrors
  \item Distinguish between real and virtual images
  \item Describe how images are formed by concave
  and convex mirrors at various object positions
  \item Apply the mirror formula and magnification
  equation to solve problems
  \item Draw ray diagrams for concave and convex
  mirrors
  \item Describe practical applications of curved
  mirrors
  \item Explain total internal reflection and
  critical angle in the context of reflection
  from mirror surfaces
\end{enumerate}
\end{objectivebox}

% ============================================================
\section{Intuition: What is Reflection?}
% ============================================================

You wake up in the morning and look in the bathroom
mirror. The face looking back at you is your own ---
but something is subtly wrong. Your right hand appears
to be the left hand of the person in the mirror. The
image is \textbf{laterally inverted}. Yet your head
is at the top and your feet at the bottom --- the image
is not inverted vertically. Why does a mirror reverse
left and right but not up and down?

The answer is that a mirror does neither. What it
actually does is reverse \textbf{front and back} ---
it swaps the direction you are facing. Your nose
points toward the mirror; the image's nose points
away from the mirror. The apparent left-right
reversal is a consequence of how we interpret ``left''
and ``right'' when comparing ourselves to our mirror
image.

This deceptively simple question --- answered by
thinking carefully about what reflection actually
does --- is the kind of thinking that Physics trains.
The laws of reflection, which we now derive, make all
of this precise.

% ============================================================
\section{The Laws of Reflection}
% ============================================================

\begin{lawbox}[Laws of Reflection]
\begin{enumerate}[noitemsep]
  \item The incident ray, the reflected ray, and the
  normal to the reflecting surface at the point of
  incidence all lie in the \textbf{same plane}.
  \item The angle of incidence equals the angle of
  reflection:
  \[
    \theta_i = \theta_r
  \]
\end{enumerate}
Both angles are measured from the \textbf{normal}
(a line perpendicular to the surface at the point
of incidence), \textit{not} from the surface itself.
\end{lawbox}

\begin{diagbox}[Laws of Reflection at a Plane Surface]
\begin{center}
\begin{tikzpicture}[scale=1.1, font=\small]

  % Mirror surface
  \fill[gray!30] (0,-0.2) rectangle (8,0);
  \draw[very thick] (0,0) -- (8,0);
  % Hatching below mirror
  \foreach \x in {0,0.4,...,7.6} {
    \draw[gray!50, thin] (\x,0) -- (\x-0.3,-0.4);
  }

  % Normal
  \draw[dashed, gray, thick] (4,-0.5) -- (4,3.5);
  \node[right, gray] at (4.1,3.3) {Normal};

  % Point of incidence
  \fill[black] (4,0) circle (3pt);

  % Incident ray
  \draw[->, very thick, codexblue, line width=2pt]
    (1.2,3.0) -- (4,0);
  \node[left, codexblue] at (1.0,3.0)
    {Incident ray};

  % Reflected ray
  \draw[->, very thick, codexred, line width=2pt]
    (4,0) -- (6.8,3.0);
  \node[right, codexred] at (7.0,3.0)
    {Reflected ray};

  % Angle of incidence arc and label (measured from the normal)
  \draw[codexblue, thick]
    (4,0.8) arc(90:133:0.8);
  \node[codexblue] at (3.42,1.15) {$\theta_i$};

  % Angle of reflection arc and label
  \draw[codexred, thick]
    (4,0.8) arc(90:47:0.8);
  \node[codexred] at (4.58,1.15) {$\theta_r$};

  % Equality annotation
  \node[codexgold, font=\bfseries] at (4,-1.0)
    {$\theta_i = \theta_r$ \quad (Law of Reflection)};

  % Wavefronts on incident beam (optional visual)
  \foreach \k in {0.5,1.0,1.5} {
    \draw[codexblue!40, thin]
      ({1.2+\k*0.7},{3.0-\k*1.0})
      -- ({1.2+\k*0.7+0.5},{3.0-\k*1.0-0.7});
  }

\end{tikzpicture}
\end{center}
\end{diagbox}

\begin{misconceptionbox}
The angle of incidence is measured from the
\textbf{normal}, NOT from the surface.
A ray travelling along the normal has
$\theta_i = 0°$ and reflects straight back
($\theta_r = 0°$). A ray grazing the surface
has $\theta_i = 90°$ and reflects at $90°$.
If you measure from the surface, you get the
\textit{complement} of the correct angle ---
answers will be wrong.
\end{misconceptionbox}

% ============================================================
\section{Reflection in Plane Mirrors}
% ============================================================

\begin{processbox}[Image Formation in a Plane Mirror]
To find where a plane mirror forms an image:
\begin{enumerate}[noitemsep]
  \item Draw two rays from the object to the mirror.
  \item Apply the law of reflection to each ray
  ($\theta_i = \theta_r$).
  \item Extend the reflected rays \textit{behind} the
  mirror (as dashed lines). Where they intersect is
  the image position.
\end{enumerate}
\end{processbox}

\begin{diagbox}[Image in a Plane Mirror]
\begin{center}
\begin{tikzpicture}[scale=1.0, font=\small]

  % Mirror
  \fill[gray!25] (5,-0.3) rectangle (5.4,5.5);
  \draw[very thick] (5,0) -- (5,5);
  \foreach \y in {0,0.4,...,5.2} {
    \draw[gray!50, thin] (5,\y) -- (5.4,\y-0.3);
  }
  \node[right, gray] at (5.5,5.3) {Mirror};

  % Object (arrow)
  \draw[->, very thick, codexblue, line width=2pt]
    (2,0) -- (2,2.5);
  \node[left, codexblue] at (1.9,1.25) {Object};
  \node[below, codexblue] at (2,0) {O};

  % Image (dashed arrow, behind mirror)
  \draw[->, very thick, codexred!50, dashed, line width=2pt]
    (8,0) -- (8,2.5);
  \node[right, codexred!70, align=center] at (8.1,1.25) {Image\\(virtual)};

  % Ray 1: incident and reflected paths on the object side.
  \draw[->, thick, codexblue] (2,2.5) -- (5,2.0);
  \draw[->, thick, codexred] (5,2.0) -- (2,1.5);
  % Dashed backward extension locates the virtual image.
  \draw[dashed, codexred!50] (5,2.0) -- (8,2.5);

  % Ray 2: a second incident ray reflects back to the left.
  \draw[->, thick, codexblue] (2,2.5) -- (5,1.5);
  \draw[->, thick, codexred] (5,1.5) -- (2,0.5);
  \draw[dashed, codexred!50] (5,1.5) -- (8,2.5);

  % Distance markers
  \draw[<->, gray, thick] (2,-0.8) -- (5,-0.8);
  \node[below, gray, font=\small] at (3.5,-0.9) {$u$};
  \draw[<->, gray, thick] (5,-0.8) -- (8,-0.8);
  \node[below, gray, font=\small] at (6.5,-0.9) {$v$};

  % Equal distance annotation
  \node[codexgold, font=\bfseries] at (5,-2.2)
    {$u = v$ \quad (image as far behind as object is in front)};

  % Apparent intersection of the backward extensions.
  \fill[codexred!40] (8,2.5) circle (4pt);
  \node[above, codexred, font=\tiny, align=center] at (8,2.55)
    {image of\\object tip};

\end{tikzpicture}
\end{center}
\end{diagbox}

\begin{defbox}[Properties of Image in a Plane Mirror]
\begin{enumerate}[noitemsep]
  \item \textbf{Same size} as the object
  ($|$image height$|$ = $|$object height$|$)
  \item \textbf{Same distance behind} the mirror
  as the object is in front
  \item \textbf{Virtual} --- the image cannot be
  formed on a screen; the reflected rays appear to
  diverge from a point behind the mirror but do not
  actually pass through that point
  \item \textbf{Erect} (right way up)
  \item \textbf{Laterally inverted} (left-right
  reversed relative to the object)
\end{enumerate}
\end{defbox}

\begin{definitionbox}[Real and Virtual Images]
A \textbf{real image} is formed where reflected
(or refracted) rays \textit{actually converge}. It
can be projected on a screen. Examples: image on a
cinema screen, image in a concave mirror when object
is beyond the focal point.

A \textbf{virtual image} is formed where reflected
(or refracted) rays \textit{appear to diverge from}
when extended backwards. The rays do not actually
pass through the image point. It cannot be projected
on a screen. Examples: image in a plane mirror,
image in a convex mirror, image in a concave mirror
when object is within the focal length.
\end{definitionbox}

% ============================================================
\section{Curved Mirrors}
% ============================================================

\begin{defbox}[Terminology for Curved Mirrors]
\begin{itemize}[noitemsep]
  \item \textbf{Centre of curvature (C):} the centre
  of the sphere of which the mirror is a part.
  \item \textbf{Radius of curvature (R):} the radius
  of that sphere.
  \item \textbf{Principal axis:} the straight line
  passing through the centre of curvature and the
  pole (centre) of the mirror.
  \item \textbf{Focal point (F):} the point where
  rays parallel to the principal axis converge after
  reflection (concave) or appear to diverge from
  after reflection (convex).
  \item \textbf{Focal length (f):} the distance from
  the pole to the focal point. $f = R/2$.
  \item \textbf{Concave mirror:} reflecting surface
  curves inward (like the inside of a spoon). Converging.
  \item \textbf{Convex mirror:} reflecting surface
  curves outward (like the outside of a spoon).
  Diverging.
\end{itemize}
\end{defbox}

\begin{diagbox}[Concave and Convex Mirror Geometry]
\begin{center}
\begin{tikzpicture}[scale=0.9, font=\small]
  % The pole P is the vertex of each mirror. C and F lie on the same
  % principal axis, with |PC|=R and |PF|=f=R/2. Curves are spherical arcs.
  \begin{scope}[shift={(0,0)}]
    \node[font=\bfseries, codexblue] at (-2,4.15)
      {Concave Mirror};

    % Concave reflecting surface faces the incident rays on the left.
    \draw[very thick, codexblue]
      plot[domain=0.5:3.5,samples=40]
      ({-4+sqrt(16-(\x-2)^2)},\x);
    \foreach \y in {0.7,1.1,1.5,1.9,2.3,2.7,3.1,3.4} {
      \pgfmathsetmacro{\xh}{-4+sqrt(16-(\y-2)^2)}
      \draw[gray!50, thin] (\xh,\y) -- ({\xh+0.22},\y+0.08);
    }

    \draw[dashed, gray] (-5.0,2) -- (1.5,2);
    \fill[codexred] (-4,2) circle (3pt);
    \node[above, codexred] at (-4,2.12) {$C$};
    \fill[codexgold] (-2,2) circle (3pt);
    \node[above, codexgold] at (-2,2.12) {$F$};
    \fill[black] (0,2) circle (3pt);
    \node[below, font=\tiny] at (0,1.88) {$P$};

    \draw[<->, gray] (0,1.2) -- (-2,1.2);
    \node[below, gray, font=\tiny] at (-1,1.12) {$f$};
    \draw[<->, gray] (0,0.75) -- (-4,0.75);
    \node[below, gray, font=\tiny] at (-2,0.65) {$R=2f$};

    % Two paraxial parallel rays strike the curved face and reflect through F.
    \foreach \y in {1.4,2.6} {
      \pgfmathsetmacro{\xh}{-4+sqrt(16-(\y-2)^2)}
      \draw[->, thick, codexblue!60] (-4.8,\y) -- (\xh,\y);
      \draw[->, thick, codexblue] (\xh,\y) -- (-2,2);
    }
  \end{scope}

  \begin{scope}[shift={(8.2,0)}]
    \node[font=\bfseries, codexred] at (2.1,4.15)
      {Convex Mirror};

    % Convex reflecting surface bulges toward the incident rays on the left.
    \draw[very thick, codexred]
      plot[domain=0.5:3.5,samples=40]
      ({4-sqrt(16-(\x-2)^2)},\x);
    \foreach \y in {0.7,1.1,1.5,1.9,2.3,2.7,3.1,3.4} {
      \pgfmathsetmacro{\xh}{4-sqrt(16-(\y-2)^2)}
      \draw[gray!50, thin] (\xh,\y) -- ({\xh+0.22},\y+0.08);
    }

    \draw[dashed, gray] (-2,2) -- (5.2,2);
    \fill[black] (0,2) circle (3pt);
    \node[below, font=\tiny] at (0,1.88) {$P$};
    \fill[codexgold!70] (2,2) circle (3pt);
    \node[below, codexgold!80, font=\tiny] at (2,1.88) {$F$ (virtual)};
    \fill[codexred!55] (4,2) circle (3pt);
    \node[above left, codexred!70, align=right, font=\tiny]
      at (4,2.12) {$C$\\(virtual)};

    % Parallel incident rays reflect to the left and diverge. Their backward
    % extensions behind the mirror meet at the virtual paraxial focus F.
    \foreach \y/\yr in {1.4/0.83,2.6/3.17} {
      \pgfmathsetmacro{\xh}{4-sqrt(16-(\y-2)^2)}
      \draw[->, thick, codexred!55] (-1.8,\y) -- (\xh,\y);
      \draw[->, thick, codexred] (\xh,\y) -- (-1.8,\yr);
      \draw[dashed, codexred!50] (\xh,\y) -- (2,2);
    }
    \node[codexred, font=\tiny, align=center] at (-1.2,3.55)
      {reflected rays\\diverge};
  \end{scope}

\end{tikzpicture}
\end{center}
\end{diagbox}

% ============================================================
\section{Ray Diagrams for Curved Mirrors}
% ============================================================

\begin{processbox}[The Three Principal Rays for Curved Mirrors]
To draw a ray diagram, use any two of these three
special rays:

\textbf{Ray 1 (Parallel ray):} A ray travelling
parallel to the principal axis reflects through the
focal point F (concave) or appears to come from F
(convex).

\textbf{Ray 2 (Focal ray):} A ray passing through
F (concave) or heading toward F (convex) reflects
parallel to the principal axis.

\textbf{Ray 3 (Centre ray):} A ray passing through
the centre of curvature C reflects back along itself
(strikes the mirror perpendicularly).
\end{processbox}

\begin{diagbox}[Concave Mirror: Object Beyond C (Real, Inverted, Diminished)]
\begin{center}
\begin{tikzpicture}[scale=0.9, font=\small]
  % Pole P=(0,3.5), centre C=(-4,3.5), and focus F=(-2,3.5).
  % The object is beyond C; the image is between C and F.
  \draw[very thick, codexblue]
    plot[domain=0.5:6.5,samples=50]
    ({-4+sqrt(16-(\x-3.5)^2)},\x);
  \foreach \y in {0.7,1.1,1.5,1.9,2.3,2.7,3.1,3.5,3.9,4.3,4.7,5.1,5.5,5.9,6.3} {
    \pgfmathsetmacro{\xh}{-4+sqrt(16-(\y-3.5)^2)}
    \draw[gray!50,thin] (\xh,\y) -- ({\xh+0.2},\y+0.08);
  }

  \draw[dashed,gray] (-7,3.5) -- (1,3.5);
  \fill[codexred] (-4,3.5) circle (3pt);
  \node[below,codexred,font=\tiny] at (-4,3.38) {$C$};
  \fill[codexgold] (-2,3.5) circle (3pt);
  \node[below,codexgold,font=\tiny] at (-2,3.38) {$F$};
  \fill[black] (0,3.5) circle (3pt);
  \node[below,font=\tiny] at (0,3.38) {$P$};

  % Object beyond C and its diminished, inverted real image.
  \draw[->,very thick,codexblue,line width=1.6pt] (-6,3.5) -- (-6,4.3);
  \node[left,codexblue,font=\tiny] at (-6.1,4.05) {Object};
  \draw[->,very thick,codexgreen,line width=1.6pt] (-2.95,3.5) -- (-2.95,3.1);
  \fill[codexgreen] (-2.95,3.1) circle (3pt);
  \node[below,codexgreen,fill=white,inner sep=1pt,font=\tiny]
    at (-2.95,2.7) {Image (real, inverted)};

  % Ray 1: parallel to the axis, then reflected through F.
  \pgfmathsetmacro{\xrayone}{-4+sqrt(16-(4.3-3.5)^2)}
  \draw[->,thick,codexblue] (-6,4.3) -- (\xrayone,4.3);
  \draw[->,thick,codexred] (\xrayone,4.3) -- (-3.5,2.875);

  % Ray 2: directed through F, then reflected parallel to the axis.
  \pgfmathsetmacro{\xraytwo}{-4+sqrt(16-(3.1-3.5)^2)}
  \draw[->,thick,codexblue] (-6,4.3) -- (\xraytwo,3.1);
  \draw[->,thick,codexred] (\xraytwo,3.1) -- (-4.8,3.1);

\end{tikzpicture}
\end{center}
\end{diagbox}

\begin{table}[H]
\centering
\caption{Summary of Images Formed by a Concave Mirror}
\label{tab:concave_images}
\small
\begin{tabular}{p{3cm}p{2cm}p{2cm}p{2cm}p{2cm}}
\toprule
\textbf{Object position} & \textbf{Image position} &
\textbf{Size} & \textbf{Nature} & \textbf{Orientation}\\
\midrule
At infinity & At F & Very small & Real & Inverted\\
Beyond C & Between F and C & Smaller & Real & Inverted\\
At C & At C & Same size & Real & Inverted\\
Between F and C & Beyond C & Larger & Real & Inverted\\
At F & At infinity & --- & --- & ---\\
Between F and P & Behind mirror & Larger & Virtual & Erect\\
\bottomrule
\end{tabular}
\end{table}

\begin{table}[H]
\centering
\caption{Summary of Images Formed by a Convex Mirror}
\label{tab:convex_images}
\small
\begin{tabular}{p{3.5cm}p{3.5cm}p{2cm}p{2cm}p{2cm}}
\toprule
\textbf{Object position} & \textbf{Image position} &
\textbf{Size} & \textbf{Nature} & \textbf{Orientation}\\
\midrule
Any position in front & Behind mirror, between P and F &
  Smaller & Virtual & Erect\\
\bottomrule
\end{tabular}
\end{table}

% ============================================================
\section{The Mirror Formula and Magnification}
% ============================================================

\begin{lawbox}[The Mirror Formula]
For a curved mirror with focal length $f$, object
distance $u$, and image distance $v$:
\[
  \frac{1}{f} = \frac{1}{u} + \frac{1}{v}
\]
\textbf{Sign convention (Real is Positive):}
\begin{itemize}[noitemsep]
  \item Distances measured from the pole P.
  \item Distances in the direction of the incident
  light are \textit{positive} (object distance $u$
  is always positive for real objects).
  \item For a concave mirror: $f$ is positive.
  \item For a convex mirror: $f$ is negative.
  \item Real images have positive $v$.
  \item Virtual images have negative $v$.
\end{itemize}
\end{lawbox}

\begin{lawbox}[Linear Magnification]
The \textbf{linear magnification} $m$ of an image
is:
\[
  m = \frac{\text{image height}}{\text{object height}}
  = \frac{h_i}{h_o} = -\frac{v}{u}
\]
\begin{itemize}[noitemsep]
  \item $|m| > 1$: image is magnified (larger than
  object)
  \item $|m| < 1$: image is diminished (smaller
  than object)
  \item $m > 0$: image is erect
  \item $m < 0$: image is inverted
\end{itemize}
\end{lawbox}

\begin{examplebox}[19.1]
\stepcounter{workedexample}
An object $4.0\ \text{cm}$ tall is placed
$30\ \text{cm}$ in front of a concave mirror
of focal length $10\ \text{cm}$.\\
(a) Find the image distance using the mirror formula.\\
(b) Find the magnification.\\
(c) Find the height of the image.\\
(d) State the nature of the image.
\end{examplebox}

\begin{solutionbox}
$u = 30\ \text{cm}$, $f = +10\ \text{cm}$ (concave)

\textbf{(a)} Mirror formula:
\[
  \frac{1}{v} = \frac{1}{f} - \frac{1}{u}
  = \frac{1}{10} - \frac{1}{30}
  = \frac{3-1}{30} = \frac{2}{30}
\]
\[
  v = \frac{30}{2} = +15\ \text{cm}
\]
Positive $v$ confirms a real image.

\textbf{(b)} Magnification:
\[
  m = -\frac{v}{u} = -\frac{15}{30} = -0.50
\]

\textbf{(c)} Image height:
\[
  h_i = m \times h_o = -0.50 \times 4.0
  = -2.0\ \text{cm}
\]
The negative sign confirms the image is inverted.

\textbf{(d)} The image is:
\textbf{real} (positive $v$),
\textbf{inverted} (negative $m$),
\textbf{diminished} ($|m| < 1$),
at $15\ \text{cm}$ in front of the mirror.
\end{solutionbox}

\begin{examplebox}[19.2]
\stepcounter{workedexample}
An object is placed $8.0\ \text{cm}$ in front of a
concave mirror. The image is virtual and is
$24\ \text{cm}$ behind the mirror.\\
(a) Calculate the focal length of the mirror.\\
(b) Calculate the radius of curvature.\\
(c) Calculate the magnification.
\end{examplebox}

\begin{solutionbox}
$u = +8.0\ \text{cm}$, $v = -24\ \text{cm}$
(negative because virtual, behind mirror)

\textbf{(a)} Focal length:
\[
  \frac{1}{f} = \frac{1}{u} + \frac{1}{v}
  = \frac{1}{8} + \frac{1}{-24}
  = \frac{3}{24} - \frac{1}{24} = \frac{2}{24}
\]
\[
  f = \frac{24}{2} = +12\ \text{cm}
\]
Positive $f$ confirms it is a concave mirror.

\textbf{(b)} Radius of curvature:
$R = 2f = 24\ \text{cm}$

\textbf{(c)} Magnification:
\[
  m = -\frac{v}{u} = -\frac{-24}{8} = +3.0
\]
Positive and greater than 1: the image is erect
and magnified (three times the object size).
This is the principle of a shaving/make-up mirror:
object placed between F and P gives a magnified,
erect, virtual image.
\end{solutionbox}

\begin{examplebox}[19.3]
\stepcounter{workedexample}
A convex driving mirror has a focal length of
$-25\ \text{cm}$. A car $3.0\ \text{m}$ behind
the mirror is being viewed.\\
(a) Find the image distance.\\
(b) Find the magnification.\\
(c) Explain why convex mirrors are used as rear-view
and security mirrors rather than plane or concave mirrors.
\end{examplebox}

\begin{solutionbox}
$f = -25\ \text{cm}$, $u = +300\ \text{cm}$

\textbf{(a)}
\[
  \frac{1}{v} = \frac{1}{f} - \frac{1}{u}
  = \frac{1}{-25} - \frac{1}{300}
  = -\frac{12}{300} - \frac{1}{300}
  = -\frac{13}{300}
\]
\[
  v = -\frac{300}{13} = -23.1\ \text{cm}
\]
Negative: image is virtual, $23.1\ \text{cm}$
behind the mirror.

\textbf{(b)}
\[
  m = -\frac{v}{u} = -\frac{-23.1}{300}
  = +0.077
\]
The image is diminished to about $7.7\%$ of the
object size.

\textbf{(c)} A convex mirror always produces an
\textit{erect, virtual, diminished} image regardless
of object distance. This gives:
\begin{itemize}[noitemsep]
  \item A wider field of view (the outward curvature
  collects light from a larger angle) --- you see more
  of the scene behind you
  \item The image is always the right way up (erect)
  \item It is always in front of the driver, behind
  the mirror surface (virtual and safe)
\end{itemize}
A plane mirror would give a narrower field of view.
A concave mirror would give a distorted, sometimes
inverted image that changes as objects move closer.
\end{solutionbox}

% ============================================================
\section{Applications of Curved Mirrors}
% ============================================================

\begin{table}[H]
\centering
\caption{Applications of Curved Mirrors}
\label{tab:mirror_applications}
\begin{tabular}{lll}
\toprule
\textbf{Mirror type} & \textbf{Application} &
\textbf{Reason}\\
\midrule
Concave & Shaving/make-up mirror &
  Magnified, erect virtual image\\
Concave & Dentist's mirror &
  Magnified, erect virtual image\\
Concave & Parabolic torch/headlight &
  Produces parallel beam from bulb at F\\
Concave & Solar furnace &
  Concentrates sunlight at F\\
Concave & Reflecting telescope &
  Collects and focuses distant light\\
Convex & Car rear-view mirror &
  Wide field of view, always erect\\
Convex & Security/shop mirror &
  Wide field of view, sees whole room\\
Convex & Road junction mirrors &
  Wide field of view for safety\\
\bottomrule
\end{tabular}
\end{table}

\begin{realworldbox}[Solar Concentrators in Africa]
Parabolic concave mirrors concentrate solar radiation
at their focal point, generating very high temperatures.
The Solar Energy Research Laboratory at Obafemi Awolowo
University (OAU) in Ile-Ife has studied parabolic
trough concentrators for water heating and cooking
applications. A parabolic solar cooker can focus enough
sunlight to boil water without any fuel --- directly
relevant to rural communities where kerosene and LPG
are expensive. The physics is simply the concave
mirror bringing all parallel solar rays to a single
focus.
\end{realworldbox}

% ============================================================
\section{Lateral Inversion and Multiple Reflections}
% ============================================================

\begin{processbox}[Image in Two Mirrors at 90°]
When two plane mirrors are placed at $90°$ to each
other, three images of an object are formed. The
third image (formed by two reflections) is not
laterally inverted --- it shows the object as others
see it, not as you see yourself in a single mirror.
This is why text appears readable in the corner of
two perpendicular mirrors, unlike in a single mirror.

For two mirrors at angle $\theta$, the number of
images formed is:
\[
  n = \frac{360°}{\theta} - 1
  \quad \text{(when } 360°/\theta \text{ is an integer)}
\]
At $90°$: $n = 360/90 - 1 = 3$ images.
At $60°$: $n = 360/60 - 1 = 5$ images.
At $45°$: $n = 360/45 - 1 = 7$ images.
\end{processbox}

% ============================================================
\section{Chapter Summary}
% ============================================================

\begin{summarybox}
\textbf{Laws of reflection:} (1) incident ray,
reflected ray, normal all in same plane;
(2) $\theta_i = \theta_r$ (angles from normal).

\textbf{Plane mirror image:} virtual, erect, same size,
laterally inverted, same distance behind as object
is in front.

\textbf{Real image:} rays actually converge;
can be projected on screen.
\textbf{Virtual image:} rays appear to diverge from;
cannot be projected.

\textbf{Concave mirror} (converging): focal length $f > 0$.
\textbf{Convex mirror} (diverging): $f < 0$; always
gives virtual, erect, diminished image.

\textbf{Mirror formula:}
$\frac{1}{f} = \frac{1}{u} + \frac{1}{v}$

\textbf{Magnification:}
$m = h_i/h_o = -v/u$

$m < 0$: inverted. $m > 0$: erect.
$|m| > 1$: magnified. $|m| < 1$: diminished.

$R = 2f$.

\textbf{Applications:}
Concave: shaving mirror, torch, telescope, solar
furnace. Convex: driving mirror, security mirror.

\textbf{Multiple mirrors:}
$n = 360°/\theta - 1$ images formed by two mirrors
at angle $\theta$.
\end{summarybox}

% ============================================================
\newpage
\begin{exercisebox}[19]

\subsection*{Section A: Multiple Choice}

\textbf{A1.} The angle of reflection is measured
from:

\mcoption{A}{The reflecting surface}
\mcoption{B}{The normal to the surface}
\mcoption{C}{The incident ray}
\mcoption{D}{The tangent to the surface}
\ansref{Ch19-A1}

\vspace{0.4cm}

\textbf{A2.} A plane mirror forms an image that is:

\mcoption{A}{Real, inverted, same size}
\mcoption{B}{Virtual, erect, same size}
\mcoption{C}{Real, erect, magnified}
\mcoption{D}{Virtual, inverted, diminished}
\ansref{Ch19-A2}

\vspace{0.4cm}

\textbf{A3.} An object is placed $20\ \text{cm}$
from a concave mirror of focal length $15\ \text{cm}$.
The image distance is:

\mcoption{A}{$-60\ \text{cm}$}
\mcoption{B}{$+60\ \text{cm}$}
\mcoption{C}{$+8.6\ \text{cm}$}
\mcoption{D}{$-8.6\ \text{cm}$}
\ansref{Ch19-A3}

\vspace{0.4cm}

\textbf{A4.} A convex mirror always produces an
image that is:

\mcoption{A}{Real, inverted, diminished}
\mcoption{B}{Real, erect, magnified}
\mcoption{C}{Virtual, erect, diminished}
\mcoption{D}{Virtual, inverted, same size}
\ansref{Ch19-A4}

\vspace{0.4cm}

\textbf{A5.} A concave mirror of radius of curvature
$40\ \text{cm}$ has a focal length of:

\mcoption{A}{$80\ \text{cm}$}
\mcoption{B}{$40\ \text{cm}$}
\mcoption{C}{$20\ \text{cm}$}
\mcoption{D}{$10\ \text{cm}$}
\ansref{Ch19-A5}

\vspace{0.4cm}

\textbf{A6.} An object is placed at the focal point
of a concave mirror. The image is:

\mcoption{A}{At the centre of curvature}
\mcoption{B}{At the focal point}
\mcoption{C}{At infinity}
\mcoption{D}{Behind the mirror}
\ansref{Ch19-A6}

\vspace{0.4cm}

\textbf{A7.} The magnification produced by a mirror
is $-3$. This means the image is:

\mcoption{A}{Erect and 3 times larger}
\mcoption{B}{Inverted and 3 times larger}
\mcoption{C}{Erect and 3 times smaller}
\mcoption{D}{Inverted and 3 times smaller}
\ansref{Ch19-A7}

\vspace{0.4cm}

\textbf{A8.} Two plane mirrors are placed at
$60°$ to each other. The number of images of an
object placed between them is:

\mcoption{A}{$3$}
\mcoption{B}{$4$}
\mcoption{C}{$5$}
\mcoption{D}{$6$}
\ansref{Ch19-A8}

\vspace{0.4cm}

\textbf{A9.} The focal length of a concave mirror
is $12\ \text{cm}$. An object is placed $6\ \text{cm}$
from the mirror. The image is:

\mcoption{A}{Real, inverted, at $12\ \text{cm}$}
\mcoption{B}{Virtual, erect, at $12\ \text{cm}$
behind mirror}
\mcoption{C}{Real, inverted, at $4\ \text{cm}$}
\mcoption{D}{Virtual, erect, at $4\ \text{cm}$}
\ansref{Ch19-A9}

\vspace{0.4cm}

\textbf{A10.} Why is a convex mirror used as a
driving mirror?

\mcoption{A}{It produces magnified images}
\mcoption{B}{It produces real images}
\mcoption{C}{It provides a wide field of view and
always erect images}
\mcoption{D}{It shows objects at their true distance}
\ansref{Ch19-A10}

\vspace{0.6cm}

\subsection*{Section B: Theory and Calculations}

\textbf{B1.} A concave mirror has a focal length
of $20\ \text{cm}$. For each of the following object
positions, use the mirror formula to find the image
position, magnification, and state the nature
(real/virtual, erect/inverted, enlarged/diminished):\\[4pt]
(a) Object at $60\ \text{cm}$\\
(b) Object at $40\ \text{cm}$\\
(c) Object at $30\ \text{cm}$\\
(d) Object at $10\ \text{cm}$\\
(e) Verify your results by sketching a ray diagram
for case (c).
\hfill\ansref{Ch19-B1}

\vspace{0.4cm}

\textbf{B2.} A security camera uses a convex mirror
of focal length $-30\ \text{cm}$ to monitor a shop.
A thief stands $2.0\ \text{m}$ from the mirror.\\[4pt]
(a) Calculate the image distance.\\
(b) Calculate the magnification.\\
(c) If the thief is $1.8\ \text{m}$ tall, how tall
does the image appear in the mirror?\\
(d) Explain why the image in a convex mirror
appears to be further away than the object
actually is. (Consider the magnification and how
we judge distance from apparent size.)\\
(e) Calculate the field of view (the maximum
angle at which an object can be seen in the mirror)
if the mirror has a diameter of $30\ \text{cm}$.
Assume the observer's eye is $50\ \text{cm}$ directly
in front of the mirror. (Use $\tan\alpha = r/d$)
\hfill\ansref{Ch19-B2}

\vspace{0.4cm}

\textbf{B3.} A dentist's concave mirror produces a
virtual image of a tooth that is $3$ times the
actual size of the tooth.\\[4pt]
(a) If the focal length of the mirror is
$2.4\ \text{cm}$, at what distance is the tooth
from the mirror?\\
(b) Find the image distance and verify that
it is virtual.\\
(c) Explain why the image must be virtual for
the dentist's application.\\
(d) If a candle is placed at $3.0\ \text{cm}$
from the same mirror, describe fully the image
formed. Is it suitable for a dentist?
\hfill\ansref{Ch19-B3}

\vspace{0.4cm}

\textbf{B4.} A torch uses a concave parabolic
reflector to produce a parallel beam of light.
The bulb filament is placed at the focal point,
$3.5\ \text{cm}$ from the vertex of the reflector.
The reflector has a diameter of $8.0\ \text{cm}$.\\[4pt]
(a) State the image position when the object
(filament) is at the focal point.\\
(b) Using $1/f = 1/u + 1/v$, verify your answer
in (a) by considering what happens as $u \to f$.\\
(c) Calculate the radius of curvature of the
reflector.\\
(d) Two identical torches are shone from the same
side: one has its filament exactly at the focal
point; the other has its filament $0.5\ \text{cm}$
beyond the focal point (further from the mirror).
Describe the difference in the beams produced.\\
(e) Explain why a parabolic reflector is used
rather than a spherical reflector for producing
a parallel beam.
\hfill\ansref{Ch19-B4}

\vspace{0.4cm}

\textbf{B5.} A student sets up an experiment to
find the focal length of a concave mirror by
measuring image distances for various object
distances.

\begin{center}
\begin{tabular}{cccc}
\toprule
$u$ (cm) & $v$ (cm) & $1/u$ (cm$^{-1}$) &
$1/v$ (cm$^{-1}$)\\
\midrule
15.0 & 60.0 & & \\
20.0 & 30.0 & & \\
25.0 & 21.4 & & \\
30.0 & 18.5 & & \\
40.0 & 15.4 & & \\
\bottomrule
\end{tabular}
\end{center}

(a) Complete the table by calculating $1/u$ and
$1/v$ for each row (to 3 significant figures).\\
(b) Plot a graph of $1/v$ against $1/u$ and draw
the best fit straight line.\\
(c) Determine the focal length from:\\
\quad (i) The gradient of the line (should be $-1$
for perfect data)\\
\quad (ii) The $y$-intercept (equals $1/f$)\\
\quad (iii) The $x$-intercept (equals $1/f$)\\
(d) Calculate the mean focal length from the three
values in (c).\\
(e) Explain one source of error in this experiment
and how it could be reduced.
\hfill\ansref{Ch19-B5}

\end{exercisebox}
